\begin{abstract}
In a two-round set reconciliation protocol based on an Invertible Bloom Lookup Table (IBLT), when the fixed-capacity difference table of the first round suffers a decoding failure because the difference is overloaded, the receiver must estimate a safe upper bound on the unknown difference size $d$ from the sum of squared counts of the difference table, so as to set the second-round table size adaptively. The standard IBLT, however, uses the set direction of each element directly as the accumulation sign, which couples the hypergraph-overlap bias selected by failure with the sign composition of the two-sided difference; a one-sided difference amplifies the estimation bias in the same direction, forcing existing schemes to rely on tedious offline simulation grid calibration.

This paper proposes the \emph{random-sign IBLT}. By introducing a pseudorandom sign independent of the position mapping, the design decouples failure selection from the sign composition in principle, while fully preserving the number of cell fields, their widths, and the peeling recovery capability. We prove that the estimator is exactly unbiased on any fixed failed hypergraph structure; we prove a Gaussian limit for the estimation error conditioned on failure when the difference size is large, and derive a closed-form Gaussian upper bound; we further characterize a discrete limit distribution dominated by duplicate edge pairs at small difference scales, and construct a Skellam upper-bound rule that strictly controls the conditional miss rate for all $d\ge2$ and requires no offline calibration. Five hundred million paired trials and end-to-end real protocol experiments spanning several orders of magnitude of load show that the random-sign IBLT removes direction sensitivity while strictly meeting the preset risk, and that the total two-round communication cost approaches the ideal lower bound for a known difference size.
\end{abstract}
